\section{現状の課題}

近年、日本国内ではeスポーツの市場規模やファン人口が増加しており、大会や地域イベントも開催されるようになっています。
一般社団法人日本eスポーツ協会(JESU)によると、日本国内のeスポーツ市場規模は、2021年の78.4億円から2024年には約161億円に拡大し、
約2.1倍となっています。また、2024年の市場ではイベント運営が37.3\%を占めています。\cite{jeSU2026} 

\begin{figure}[H]
  \centering
  \makeatletter
  \def\pgfpie@numbertext#1{%
    \pgfpie@ifhidenumber{}{%
      % ↓ 1未満の数値の場合はグラフ内部への数字描画を非表示にする
      \ifdim #1pt<1pt
      \else
        \pgfpie@beforenumber#1\pgfpie@afternumber
      \fi
    }%
  }
  \makeatother
  \begin{tikzpicture}
    \pie[
      radius=3,
      sum=auto,
      text=pin,
      explode={0,0,0,0,0,0,0.3,0},
      color={orange!80, red!60, olive!70, lime!60!olive, violet!50, blue!40, cyan!70!blue, red!70!black}
    ]{
      9.7/{放映権},
      36.3/{スポンサー},
      1.9/{チケット},
      7.4/{アイテム課金・賞金},
      0.8/{グッズ(0.8)},
      0.3/{\raisebox{-1.5em}{著作権許諾(0.3)}},
      37.3/{イベント運営},
      6.2/{ストリーミング}
    }
  \end{tikzpicture}
  \caption{日本国内eスポーツ市場の割合（出典：一般社団法人日本eスポーツ協会，「eスポーツの市場と推移」）}
  \label{fig:revenue}
\end{figure}


さらに、国内eスポーツファン数は2021年の743万人から2024年には967万人に増加し、約30\%増加しています。
この推移を図\ref{fig:trend}に示します。

大会開催数についても、ゲーム大会運営プラットフォーム「Tonamel」内で開催されたゲーム大会は、2018年の年間898大会から2021年には年間10,094大会へと増加したことが報告されています[2]。
\footnote[2]{株式会社カヤック，「事業計画及び成長可能性に関する事項」，2022年3月29日．株式会社カヤック（2026年10月30日閲覧）}

\begin{figure}[htbp]
\centering

\begin{tikzpicture}
\begin{axis}[
    ybar,
    width=12cm,
    height=8cm,
    xlabel={年},
    ylabel={大会数（件）},
    title={Tonamelで開催されたゲーム大会数の推移},
    xtick={1,2,3,4,5},
    xticklabels={2017年,2018年,2019年,2020年,2021年},
    ymin=0,
    ymax=11500,
    ytick distance=2000,
    nodes near coords,
    nodes near coords style={font=\small},
    bar width=25pt,
    grid=major,
]

\addplot[fill=orange!75] coordinates {
    (1,203)
    (2,898)
    (3,1253)
    (4,5812)
    (5,10094)
};

\end{axis}
\end{tikzpicture}

\caption{Tonamelで開催されたゲーム大会数の推移}

\end{figure}



このことから、ゲーム大会を開催・運営する機会が拡大していることから大会管理のための必要性が増加していることが伺えます。 このことから、ゲーム大会を開催・運営する機会が拡大しているら大会管理のための必要性が増加していることが伺えます。

\begin{figure}[H]
  \centering
  \begin{tikzpicture}
    \begin{axis}[
      ybar,
      width=0.85\linewidth,
      height=7cm,
      bar width=1.2cm,
      symbolic x coords={2020年,2021年,2022年,2023年,2024年},
      xtick=data,
      ymin=0, ymax=10000,
      ytick={0,2000,4000,6000,8000,10000},
      scaled y ticks=false,
      yticklabel={\pgfmathprintnumber[fixed,1000 sep={,}]{\tick}},  
      title={\small （千人）},
      title style={anchor=east, at={(0,1.02)}, xshift=-0.5cm},
      nodes near coords,
      nodes near coords style={/pgf/number format/fixed, /pgf/number format/1000 sep={,}},
      ymajorgrids=true,
      grid style={dashed, gray!50},
      enlarge x limits=0.15
    ]
      \addplot[fill=violet!70!black, draw=none] coordinates {
        (2020年,6859)
        (2021年,7434)
        (2022年,7759)
        (2023年,8560)
        (2024年,9673)
      };
    \end{axis}
  \end{tikzpicture}
  \caption{日本国内eスポーツファン数の推移（出典：\cite{jeSU2026}）}
  \label{fig:trend}
\end{figure}

一方、大会運営では参加者情報や試合結果など、さまざまな情報を管理する必要があります。
「Tonamel」などの既存の大会運営プラットフォームには、参加登録やトーナメント管理などの機能がありますが、参加者への連絡にDiscordやメール、試合結果の集計にExcelやGoogleスプレッドシートなど、別のツールを併用する場合もあります。
\footnote[3]{株式会社エイプリルナイツ所属者，「180名が参加するeスポーツ大会（Apex Legends）の集計処理を爆速で処理する話」，Qiita，2022年12月10日． }
また、JCG\footnote[4]{JCGとは、日本国内でeスポーツ大会の企画・運営や配信などを行っているeスポーツ関連企業・大会運営サービス}の大会参加に関する解説でも、大会によってメール、大会サイトのチャット、Discordなど連絡方法が異なることが紹介されています。
このように複数のツールを使用すると、主催者は情報の管理や更新に手間がかかり、競技参加者も試合時間や対戦相手を確認するために、複数の場所を確認する必要があります。

また、観戦者にとっても、大会全体の進行状況や注目選手の試合情報を簡単に確認できることが重要です。
eスポーツのオンライン観戦需要に関する研究では、観戦者の心理的な動機や視聴環境など、さまざまな要因がオンライン観戦に関係することが調査されています。
\footnote[5]{Tyreal Yizhou Qian, James Jianhui Zhang, Jerred Junqi Wang, John Hulland, “Beyond the Game: Dimensions of Esports Online Spectator Demand,” Communication \& Sport, Vol. 8, No. 6, pp. 825–851, 2020．}
このことから、観戦者が大会や選手の情報を確認しやすい環境を整えることも重要であると考えられます。
特に特定の選手に注目している観戦者にとっては、その選手の試合予定や対戦相手、試合結果をまとめて確認できる仕組みが役立つと考えられます。
しかし、トーナメント表や試合結果、ライブ配信などが別々のページやサービスで提供されている場合、観戦者は必要な情報を個別に確認しなければなりません。

既存の大会運営プラットフォームには、大会作成、参加者登録、トーナメント設定、勝敗管理などの機能が提供されています。
既存のプラットフォームにおける運営の効率化だけでなく、観戦者が大会情報とライブ配信を関連付けて確認できる仕組みも重要です。
特に、小規模から中規模の大会では、専用システムの導入費用や運用費用が負担となる場合もあり、利用しやすい大会運営システムが求められています。

そこで本提案書では、大会運営に必要な情報を一元管理し、特に観戦者が大会の進行状況や注目選手の試合を追いやすくする「小規模ゲーム大会運営システム(Play Line)」を提案します。
本システムでは、観戦者向けに大会全体のトーナメント表や試合結果に加え、注目選手の過去の試合結果、現在の試合状況、次の対戦相手やトーナメント上の位置をまとめて確認できる機能を提供します。さらに、主催者向けには参加登録や試合結果の管理機能、競技参加者向けには試合予定や対戦相手の確認機能を提供します。これにより、観戦者の利便性向上を中心に、大会運営の効率化と競技参加者への情報提供も実現することを目指します。